\chapter{Introduction}

Segmentation of three-dimensional (3D) cellular images is an essential part of many research pipelines, including pathological diagnosis \cite{fischmankeratinocytes}, cell tracking for proliferation \cite{Betjes2024.10.11.617799}, and embryonic mitosis development studies \cite{Stegmaier16}. Recently emerged deep learning-based methods have shown great success in target-oriented domains \cite{fischmankeratinocytes, Zhang2024.04.05.588365}. However, they show limited generality when applied to multimodal environments \cite{lee2022mediarharmonydatacentricmodelcentric}. 

Classical approaches range from using dataset-specific heuristics for thresholding, to using shape descriptors such as Laplacian-of-Gaussian in multiple scale spaces, paired with techniques using seed points to expand into contour converging algorithms like watershed \cite{inproceedings}. 
With the success of Convolutional Neural Networks (CNNs) for various image processing tasks \cite{Krizhevsky.2012}, more robust methods developed specifically for cell analysis. \cite{he2018maskrcnn, ronneberger2015unetconvolutionalnetworksbiomedical}. To further improve robustness in dense cell regions, recent work has developed advanced latent representations as shape descriptors, that can be processed into individual instances with strong performance across various modalities. StarDist \cite{Schmidt_2018} predicts distance rays for candidate pixels with a separately predicted object probability, along with Non-Maximum Suppression postprocessing for dense instance distinction. Cellpose adapts the learning objective to a simulated heat diffusion vector field, predicting vertical and horizontal gradients that point towards cell centers. A subsequent gradient integration aided by a separately predicted per-pixel cellprobability leads to accurate contour instance segmentation.  
While these learning objectives improved generality by a large margin, new methods have arisen for intermediate feature acquisition. While initial architectures relied fully on convolutional backbones, new approaches towards transformer-based feature encoders have been introduced into the pipeline \cite{lee2022mediarharmonydatacentricmodelcentric}, as Vision Transformers \cite{dosovitskiy2021imageworth16x16words} have shown great success on various vision tasks.
While inductive bias from intrinsic feature acquisition mechanisms like the convolution operation enable CNNs to perform well with limited data, transformer-based backbones need to compensate the lack of inductive bias from pretraining with a strong visual pretext task \cite{ryali2023hierahierarchicalvisiontransformer}. This limits the potential of transformer-based backbones for 3D feature acquisition, as labeled 3D microscopy data is limited. However, slice-wise stitching and feature accumulation across non-axial directions like used in Cellpose, have shown great success with 2D CNN backbones before \cite{Stringer.2020}. 

Using this feature accumulation across different slicing directions, we try to leverage the strongly generalizable foundation models that operate in 2D, to predict robust, multi-domain cell features in 3D. As the authors of the original Cellpose released their new updated version whilst our development using this exact method of combining a strong 2D foundational model with their learning objective, they set a strong baseline for our models.

In summary, we present a new method for robust and generalizable 3D cell segmentation by adapting the Cellpose slicing flow accumulation to a transformer backbone with strong inductive bias from pretraining. We use Mediar, an established transformer-based model with strong 2D performance on cellular segmentation and extend it to 3D with the flow accumulation method of Cellpose. We will present related work in chapter \ref{State of the art} to form the basis of the techniques used in our methodology. In chapter \ref{Methods}, we will explain in detail how we set up our Mediar extension to 3D in section \ref{Mediar3D} and outline the architectural basis of the network, in order to illustrate the rationale behind certain design choices that we will take, to try to improve our model. After proposing architectural design choices, we will outline our training pipeline and discuss challenges of partially labeled datasets and feasibility of 3D feature training in section \ref{training pipeline}. To benchmark model performance depending on both architectural design choices and pretraining biases, we create 3 candidate models: First, we extend Mediar with its weights from the original paper to 3D using the Cellpose flow accumulation. Second, we take the first model and retrain it from scratch using a larger pretraining dataset. Third and last, we take the modified Mediar architecture, derived from our theoretical considerations from the methods section, and train it also on the larger pretraining dataset. This pretraining process, described in section \ref{chapter:4.1}, builds the basis for our 3 candidate models that we benchmark against the state-of-the-art. We set up experiments for testing network training with partially annotated data in section \ref{Finetuning} and explain the setup for the benchmarking experiments in section \ref{Inference}. Finally, we present our results in chapter \ref{results}.
